\chapter{Methodology}

\section{Methodology Overview}

The methodology of FireWatch is built around two complementary forms of
wildfire monitoring. The first is weather-based risk prediction, which
estimates daily fire-weather danger for the Karabük region. The second is
visual fire/smoke detection, which provides camera-based evidence when fire or
smoke appears in a monitoring feed. These two modules are methodologically
separate: the prediction model estimates fire-weather risk, while the detection
model identifies visible fire or smoke evidence.

The prediction workflow begins with weather variables collected for Karabük
from Open-Meteo, while the processed dataset includes a Fire Weather Index
(FWI) target associated with official CEMS/Copernicus fire-danger data. The
prepared dataset stores this supervised learning target in the column
\texttt{FWI}. The repository stores the processed FWI target used for
modelling, while the raw external download workflow is treated as a
data-preparation step rather than a runtime dependency. The project does not
calculate the full Canadian FWI equation chain at runtime. Instead, the FWI
value is used as the training target, and the deployed model learns to
approximate it from accessible weather variables and engineered temporal
features. This choice keeps the operational pipeline practical while anchoring
the model to an established fire-danger indicator.

The complete workflow is as follows: weather records are collected and aligned
by date, the May--October seasonal window is selected, raw weather variables
are transformed into a locked 34-feature schema, a Stage 1 regression model
predicts continuous FWI, and a Stage 2 safety classifier estimates whether the
day should be treated as high risk. The resulting decision is stored and shown
on the dashboard. In parallel, the YOLOv8 fire/smoke detector processes camera
or drone-ready video sources and creates persistent detection alerts when fire
or smoke is detected.

\begin{figure}[htbp]
  \centering
  \includegraphics[width=0.92\textwidth]{end_to_end_workflow_flowchart.png}
  \caption{End-to-end FireWatch workflow from weather input and visual monitoring to prediction, detection alerts, persistence, and dashboard review.}
  \label{fig:end-to-end-workflow}
\end{figure}

Figure~\ref{fig:end-to-end-workflow} summarizes the full methodology. It
shows that weather-based prediction and YOLOv8 visual detection are connected
inside the same system, but they answer different operational questions.

\section{Study Region and Data Sources}

The study region is Karabük, Türkiye, with the system configuration using the
coordinates \(41.2061^\circ\) latitude and \(32.6204^\circ\) longitude for
model-input weather collection. Karabük was selected as a regional pilot
because the project aims to support local wildfire preparedness rather than
produce a generic non-local risk model.

Weather variables are obtained from Open-Meteo. The data pipeline separates
historical model-input preparation from operational forecast-based inference:
historical archive data are used for past records, while forecast data can be
used when the target date falls on or after the current Istanbul-local date.
The weather variables used by the model are daily aggregate values rather than
display-only live weather values. This separation is important because the
dashboard can show current weather for operator awareness, but the prediction
model must receive inputs that match the training dataset semantics.

The supervised learning target is the FWI value stored in the dataset as
\texttt{FWI}. The final FireWatch methodology treats this target as the daily
fire-danger signal used during training. The Canadian FWI system is an
established fire-danger rating framework \cite{stocks1989}; however, FireWatch
does not reproduce the full chain of FWI sub-index equations during runtime.
Instead, the target teaches the model how daily weather and recent drying
patterns relate to FWI.

\section{Data Collection Period and Seasonal Filtering}

The processed FireWatch dataset covers seasonal records from 1 May 2012 to
31 October 2025. The processed file contains 2,576 daily rows and
uses the months May, June, July, August, September, and October. This confirms
that the training and evaluation data are restricted to the warm-season period
most relevant to wildfire risk in the project scope.

Seasonal filtering is used for methodological reasons. Including winter and
low-risk months would add many examples whose weather patterns are less
informative for fire-season decisions. This could make the model appear more
stable while weakening its focus on the period when wildfire preparedness is
operationally important. Restricting the dataset to May--October keeps the
learning problem aligned with the dashboard's purpose: supporting daily
fire-season monitoring for Karabük.

\section{Daily Weather Aggregation and Model Input Alignment}

The current model-input pipeline does not use noon-only extraction. The
Open-Meteo input pipeline requests daily aggregate variables with semantics
that match the training CSV. These daily inputs are:
\texttt{temperature\_2m\_max}, \texttt{relative\_humidity\_2m\_min},
\texttt{wind\_speed\_10m\_max}, \texttt{precipitation\_sum},
\texttt{cloud\_cover\_mean}, \texttt{shortwave\_radiation\_sum}, and
\texttt{et0\_fao\_evapotranspiration}. These variables are mapped into the
model schema as \texttt{temperature}, \texttt{rh}, \texttt{ws}, and
\texttt{precip} where needed. The
\texttt{soil\_moisture\_0\_to\_7cm\_mean} variable is added separately through
the soil-moisture resolver.

This design converts weather data into one representative daily model row. It
does not attempt to reconstruct every hourly fire-weather observation or the
full Canadian FWI state chain. The methodological reason is consistency:
Stage 1 was trained on daily model-input variables, so scheduled and manual
risk checks must use the same daily aggregate meanings. The official FWI target
is daily and is used only as a supervised training label; it is not included as
a runtime input feature.

\section{Data Cleaning and Alignment}

The prepared dataset is organised as daily records and aligned using the
\texttt{date} column as the main temporal key. Additional date-related columns,
such as \texttt{day}, \texttt{month}, and \texttt{year}, are retained to
support seasonal filtering and calendar-based feature engineering. Before model
training and inference, the records are sorted chronologically so that rolling
and temporal-memory features are computed in the correct order.

Column names are standardised to match the model feature schema. For example,
weather variables such as temperature, relative humidity, wind speed, and
precipitation are represented consistently as \texttt{temperature},
\texttt{rh}, \texttt{ws}, and \texttt{precip}. This standardisation is
important because the same feature schema must be used during both model
training and operational inference. The high-risk classification target is
derived from the official \texttt{FWI} value as \texttt{target\_ge\_35}, where
a day is considered high risk when \( \mathrm{FWI} \geq 35 \).

Target leakage is controlled during temporal feature construction. Rolling and
exponentially weighted features are computed from shifted weather series so
that they describe previous conditions rather than accidentally including the
same day's target-related information. As a result, rolling rainfall, humidity,
temperature, wind, EWMA, and dry-day features represent historical weather
memory. The official \texttt{FWI} value is used for supervised training and
evaluation, but it is not included as one of the 34 runtime input features.

\section{Feature Engineering}

FireWatch uses a locked schema of 34 input features for the prediction model.
These features are designed to approximate the environmental behaviour
represented by official FWI while relying on accessible weather variables and
locally computed temporal features. The schema contains eight raw daily weather
features and twenty-six engineered features.

The raw daily weather inputs are:
\texttt{temperature}, \texttt{rh}, \texttt{ws}, \texttt{precip},
\texttt{cloud\_cover\_mean}, \texttt{shortwave\_radiation\_sum},
\texttt{et0\_fao\_evapotranspiration}, and \\
\texttt{soil\_moisture\_0\_to\_7cm\_mean}. These variables represent immediate
fire-weather conditions such as heat, humidity stress, wind exposure, rainfall,
solar energy, evaporation demand, and near-surface soil moisture.

Seasonal and cyclical features are represented by \texttt{day\_sin},
\texttt{day\_cos}, \texttt{month\_sin}, and \texttt{month\_cos}. These
features allow the model to learn seasonal patterns without treating calendar
day or month as a simple linear value.

Derived fire-weather features include \texttt{vpd},
\texttt{fuel\_drying\_rate}, \texttt{hdw}, \texttt{wind\_squared}, and
\texttt{dew\_point}. These variables describe humidity stress, drying
potential, wind-driven spread potential, and the relationship between
temperature and atmospheric moisture. They help the model represent
fire-relevant physical behaviour without calculating the full FWI equation
chain during runtime.

Rolling history features describe recent weather memory. They include
\texttt{precip\_sum\_3d}, \texttt{precip\_sum\_7d},
\texttt{precip\_sum\_30d}, \texttt{t\_mean\_3d}, \texttt{t\_mean\_7d},
\texttt{rh\_min\_3d}, \texttt{rh\_min\_7d}, \texttt{rh\_mean\_3d},
\texttt{rh\_mean\_7d}, \texttt{ws\_mean\_3d}, \texttt{ws\_mean\_7d},
\texttt{ws\_max\_3d}, and \texttt{ws\_max\_7d}. These variables capture
accumulated rainfall memory, persistent heat, low-humidity periods, and recent
wind exposure.

Exponentially weighted memory features are \texttt{ewma\_t},
\texttt{ewma\_rh}, and \texttt{ewma\_precip}. They assign stronger weight to
recent weather while still preserving short-term historical influence. The
final dryness
persistence feature is \texttt{consecutive\_dry\_days}, which measures how
long the recent period has remained dry. Together, these feature groups allow
the model to represent short-term conditions, accumulated drying, seasonality,
and recovery after rainfall.

\begin{figure}[htbp]
  \centering
  \includegraphics[width=0.82\textwidth]{feature_groups_34.png}
  \caption{Grouped view of the 34 FireWatch prediction features used by the deployed Stage 1 model.}
  \label{fig:feature-groups}
\end{figure}

The feature groups in Figure~\ref{fig:feature-groups} keep the prediction
schema explainable. They also help ensure that training and runtime inference
use the same model inputs.

\section{Stage 1: Wildfire Risk Regression Model}

Stage 1 is the core numerical prediction model in the FireWatch pipeline. It
uses a \texttt{HistGradientBoostingRegressor} to predict the continuous FWI
target from the 34 engineered input features. The output of this model is
\texttt{predicted\_fwi}, which represents the estimated daily fire-weather
danger for the target date.

Regression is suitable for Stage 1 because FWI is a continuous intensity index
rather than a binary class label. Predicting a continuous value preserves
information about the severity of fire-weather conditions and shows how close a
day is to the high-risk boundary. This value can also be displayed on the
dashboard as part of the risk explanation for operators.

The model is implemented within a scikit-learn pipeline that handles missing
values consistently using median imputation. During training, sample weighting
is used so that mid-FWI and high-FWI cases receive greater attention than
low-risk cases. This is important because high-risk conditions are
operationally more critical, even if they occur less frequently than low-risk
days.

The stored Stage 1 metadata reports the model configuration, feature count,
training/test split, and validation metrics. These metrics are discussed later
in the Results and Discussion chapter, while this section focuses on the
modelling procedure.

\section{Stage 2: High-Risk Safety Classifier}

Stage 2 is a safety-oriented classification layer built on top of the Stage 1
regression output. It uses a \texttt{RandomForestClassifier} trained on Stage 1
out-of-fold predictions together with three support features:
\texttt{rh}, \texttt{ws}, and \texttt{fuel\_drying\_rate}. The full Stage 2
input vector is therefore \texttt{predicted\_fwi}, \texttt{rh}, \texttt{ws},
and \texttt{fuel\_drying\_rate}. Its target is \texttt{target\_ge\_35}, where
a day is labelled as high risk when \( \mathrm{FWI} \geq 35 \).

Stage 2 does not replace the Stage 1 regression model. Instead, it supports the
final operational decision by focusing on high-risk cases, especially near the
decision boundary. The regression output remains the main signal: if
\texttt{predicted\_fwi} is at least 35, the day is classified as high risk. If
\texttt{predicted\_fwi} falls in the grey zone from 28 to below 35, the Stage 2
probability can promote the day to high risk when it reaches the configured
probability threshold. The stored threshold is 0.10.

This two-stage structure is used because high FWI conditions are more important
from an operational safety perspective than average low-risk predictions. The
regression model preserves the continuous risk intensity, while the classifier
adds a decision-support layer focused on avoiding missed high-risk cases. The
stored Stage 2 metadata reports the training/test split, probability threshold,
and classification metrics. These results are discussed later in the Results
and Discussion chapter.

\begin{figure}[htbp]
  \centering
  \includegraphics[width=0.90\textwidth]{prediction_stage_flowchart.png}
  \caption{Prediction-stage workflow showing feature construction, Stage 1 FWI regression, Stage 2 high-risk classification, and final decision logic.}
  \label{fig:prediction-pipeline}
\end{figure}

Figure~\ref{fig:prediction-pipeline} highlights why the prediction stage is
not a single binary classifier. Continuous FWI estimation is preserved first,
then the high-risk safety layer supports the final operational flag.

\section{Fire and Smoke Detection Methodology}

The visual detection module is separate from the weather-based prediction
pipeline. It does not predict FWI and does not modify the wildfire risk model.
Its role is to inspect camera frames and identify visible fire or smoke
evidence that can be stored and reviewed by an operator.

The active detector is a custom YOLOv8 object detector loaded from
\texttt{best.pt} as the final fire/smoke detection model. The detector
supports two classes: fire and smoke. Runtime labels are normalised so that
Fire, fire, flame, and flames map to \texttt{fire}, while Smoke and smoke map
to \texttt{smoke}. Unknown labels are handled safely rather than being forced
into the fire class.

YOLOv8 is suitable for this module because it performs object localisation,
not only image-level classification. It can produce bounding boxes,
confidence values, and class labels on live camera frames, which are useful
for visual confirmation and alert review. In the current prototype, the
default detection confidence is configured through \texttt{MONITORING\_YOLO\_CONF}
with a default value of 0.4, and the default image size is 640.

The current detection methodology remains frame-based. False positives from
bright light, glare, reflections, or visually similar objects remain a known
limitation. More advanced robustness methods, such as multi-frame
confirmation, glare filtering, bounding-box area filtering, and extended field
validation, are left for future work.

\begin{figure}[htbp]
  \centering
  \includegraphics[width=0.90\textwidth]{detection_stage_flowchart.png}
  \caption{Detection-stage workflow for camera or drone-ready frames, YOLOv8 fire/smoke inference, alert creation, and dashboard review.}
  \label{fig:detection-pipeline}
\end{figure}

Figure~\ref{fig:detection-pipeline} shows the monitoring path as a separate
evidence pipeline. It does not write prediction runs or alter FWI estimates.

\section{Operational Decision Workflow}

At runtime, the operational prediction workflow starts by selecting the target
date in the Istanbul operational calendar. The system obtains daily
Open-Meteo model inputs, builds the 40-day history window needed for rolling
and exponentially weighted features, generates a single 34-feature row, and
runs the stacked prediction pipeline. The output includes
\texttt{predicted\_fwi}, \texttt{high\_risk\_probability}, a binary high-risk
flag, and the decision reason.

Scheduled checks are configured for 09:00, 11:00, and 15:00 Istanbul time.
Each scheduled check follows the same data-to-decision path as a manual risk
check and stores the result for dashboard review. The system distinguishes
scheduled and manual runs so that operators can audit when a decision was
created and why it was created.

The dashboard consumes the prediction output as a decision-support signal. It
shows risk information, recent runs, model/system status, and map context.
SQLite stores prediction history and detection alerts so that results do not
disappear after a page refresh or backend restart. Detection alerts preserve
the fire/smoke label and can be reviewed through the Detection Alerts
interface.

Drone-ready support is included as an operational layer rather than as a claim
of full autonomous deployment. The project supports mock and Tello-oriented
configuration, demo patrol behaviour, and camera-stream integration. Physical
drone actions remain protected by operator-controlled safety gates, while full
autonomous waypoint patrol, regulatory flight approval, and production safety
certification remain future work.

\section{Methodology Summary}

The final FireWatch methodology uses a regional, two-stage, decision-support
workflow for Karabük. The processed FWI target provides the supervised
training signal, while Open-Meteo weather variables and engineered memory
features provide practical runtime inputs. Stage 1 estimates continuous FWI,
Stage 2 adds a safety layer for high-risk decisions, and the YOLOv8 detector
supplies visual fire/smoke evidence through a separate monitoring path.
